\section{Differentiated Poisson summation}\label{app:poisson}

Put $m=n+1$ in \eqref{eq:character}.  Since
$n(n+2)=m^2-1$,
\begin{equation}\label{eq:poisson-start}
p_t(\theta)=\frac{e^{t/2}}{\sin\theta}
\sum_{m\ge1}m e^{-m^2t/2}\sin(m\theta).
\end{equation}
The Jacobi identity in the Fourier convention used here is
\[
\sum_{m\in\mathbb Z}e^{-m^2t/2}e^{im\theta}
=\sqrt{\frac{2\pi}{t}}
\sum_{k\in\mathbb Z}e^{-(\theta+2\pi k)^2/(2t)}.
\]
Normal convergence permits differentiation in $\theta$.  Multiplying the
derivative by $-1$ and taking real quantities gives
\[
2\sum_{m\ge1}m e^{-m^2t/2}\sin(m\theta)
=\sqrt{\frac{2\pi}{t}}\frac1t
\sum_{k\in\mathbb Z}(\theta+2\pi k)
e^{-(\theta+2\pi k)^2/(2t)}.
\]
Substitution in \eqref{eq:poisson-start} yields \eqref{eq:image}, including
the factor $C_t=e^{t/2}\sqrt{2\pi}/(2t^{3/2})$.

At $\theta=0$, pair $k$ and $-k$ before taking the limit.  At
$\theta=\pi$, pair the images whose arguments tend to opposite odd multiples
of $\pi$.  These operations are justified termwise by Gaussian domination.
Thus the apparent zeros of the Weyl denominator are removable in the complete
sum, but not necessarily in an unpaired truncation.
